\section{The divisor associated to an invertible sheaf}
\label{section-divisor-invertible-sheaf}

\noindent
The following definition is the analogue of
Divisors, Definition \ref{divisors-definition-divisor-invertible-sheaf}
in our current setup.

\begin{definition}
\label{definition-divisor-invertible-sheaf}
Let $(S, \delta)$ be as in Situation \ref{situation-setup}.
Let $X$ be locally of finite type over $S$. Assume $X$ is
integral and $n = \dim_\delta(X)$.
Let $\mathcal{L}$ be an invertible $\mathcal{O}_X$-module.
\begin{enumerate}
\item For any nonzero meromorphic section $s$ of $\mathcal{L}$
we define the {\it Weil divisor associated to $s$} is the
$(n - 1)$-cycle
$$
\text{div}_\mathcal{L}(s) =
\sum \text{ord}_{Z, \mathcal{L}}(s) [Z]
$$
defined in Divisors, Definition
\ref{divisors-definition-divisor-invertible-sheaf}.
This makes sense because Weil divisors have $\delta$-dimension $n - 1$
by Lemma \ref{lemma-divisor-delta-dimension}.
\item We define {\it Weil divisor associated to $\mathcal{L}$} as
$$
c_1(\mathcal{L}) \cap [X] =
\text{class of }\text{div}_\mathcal{L}(s) \in \CH_{n - 1}(X)
$$
where $s$ is any nonzero meromorphic section of $\mathcal{L}$ over
$X$. This is well defined by
Divisors, Lemma \ref{divisors-lemma-divisor-meromorphic-well-defined}.
\end{enumerate}
\end{definition}

\noindent
Let $X$ and $S$ be as in Definition \ref{definition-divisor-invertible-sheaf}
above. Set $n = \dim_\delta(X)$. It is clear from the definitions that
$Cl(X) = \CH_{n - 1}(X)$ where $Cl(X)$ is the Weil divisor class group of $X$
as defined in Divisors, Definition \ref{divisors-definition-class-group}.
The map
$$
\Pic(X) \longrightarrow \CH_{n - 1}(X), \quad
\mathcal{L} \longmapsto c_1(\mathcal{L}) \cap [X]
$$
is the same as the map $\Pic(X) \to Cl(X)$ constructed in
Divisors, Equation (\ref{divisors-equation-c1}) for arbitrary
locally Noetherian integral schemes. In particular, this map
is a homomorphism of abelian groups, it is injective if $X$ is
a normal scheme, and an isomorphism if all local rings of $X$
are UFDs. See Divisors, Lemmas \ref{divisors-lemma-normal-c1-injective} and
\ref{divisors-lemma-local-rings-UFD-c1-bijective}.
There are some cases where it is easy to compute the
Weil divisor associated to an invertible sheaf.

\begin{lemma}
\label{lemma-compute-c1}
Let $(S, \delta)$ be as in Situation \ref{situation-setup}.
Let $X$ be locally of finite type over $S$. Assume $X$ is
integral and $n = \dim_\delta(X)$.
Let $\mathcal{L}$ be an invertible $\mathcal{O}_X$-module.
Let $s \in \Gamma(X, \mathcal{L})$ be a nonzero global section.
Then
$$
\text{div}_\mathcal{L}(s) = [Z(s)]_{n - 1}
$$
in $Z_{n - 1}(X)$ and
$$
c_1(\mathcal{L}) \cap [X] = [Z(s)]_{n - 1}
$$
in $\CH_{n - 1}(X)$.
\end{lemma}

\begin{proof}
Let $Z \subset X$ be an integral closed subscheme of
$\delta$-dimension $n - 1$. Let $\xi \in Z$ be its generic
point. Choose a generator $s_\xi \in \mathcal{L}_\xi$.
Write $s = fs_\xi$ for some $f \in \mathcal{O}_{X, \xi}$.
By definition of $Z(s)$, see
Divisors, Definition \ref{divisors-definition-zero-scheme-s}
we see that $Z(s)$ is cut out by a quasi-coherent
sheaf of ideals $\mathcal{I} \subset \mathcal{O}_X$ such
that $\mathcal{I}_\xi = (f)$. Hence
$\text{length}_{\mathcal{O}_{X, \xi}}(\mathcal{O}_{Z(s), \xi})
=
\text{length}_{\mathcal{O}_{X, \xi}}(\mathcal{O}_{X, \xi}/(f))
=
\text{ord}_{\mathcal{O}_{X, \xi}}(f)$ as desired.
\end{proof}

\noindent
The following lemma will be superseded by the more general
Lemma \ref{lemma-flat-pullback-cap-c1}.

\begin{lemma}
\label{lemma-flat-pullback-divisor-invertible-sheaf}
Let $(S, \delta)$ be as in Situation \ref{situation-setup}.
Let $X$, $Y$ be locally of finite type over $S$. Assume $X$, $Y$
are integral and $n = \dim_\delta(Y)$.
Let $\mathcal{L}$ be an invertible $\mathcal{O}_Y$-module.
Let $f : X \to Y$ be a flat morphism of relative dimension $r$. Then
$$
f^*(c_1(\mathcal{L}) \cap [Y]) = c_1(f^*\mathcal{L}) \cap [X]
$$
in $\CH_{n + r - 1}(X)$.
\end{lemma}

\begin{proof}
Let $s$ be a nonzero meromorphic section of $\mathcal{L}$.
We will show that actually
$f^*\text{div}_\mathcal{L}(s) = \text{div}_{f^*\mathcal{L}}(f^*s)$
and hence the lemma holds.
To see this let $\xi \in Y$ be a point and let $s_\xi \in \mathcal{L}_\xi$
be a generator. Write $s = gs_\xi$ with $g \in R(Y)^*$.
Then there is an open neighbourhood $V \subset Y$ of $\xi$
such that $s_\xi \in \mathcal{L}(V)$ and such that $s_\xi$ generates
$\mathcal{L}|_V$. Hence we see that
$$
\text{div}_\mathcal{L}(s)|_V = \text{div}_Y(g)|_V.
$$
In exactly the same way, since $f^*s_\xi$ generates $f^*\mathcal{L}$
over $f^{-1}(V)$ and since $f^*s = g f^*s_\xi$ we also
have
$$
\text{div}_\mathcal{L}(f^*s)|_{f^{-1}(V)}
=
\text{div}_X(g)|_{f^{-1}(V)}.
$$
Thus the desired equality of cycles over $f^{-1}(V)$ follows from the
corresponding result for pullbacks of principal divisors, see
Lemma \ref{lemma-flat-pullback-principal-divisor}.
\end{proof}



